% ═══════════════════════════════════════════════════════════════════
% CHAPITRE 1 : CONTEXTE GÉNÉRAL
% ═══════════════════════════════════════════════════════════════════

\chapter{Contexte Général}

\section{Présentation de l'organisme d'accueil}

L'\textbf{ARSII} (Association de Recherche Scientifique et Innovation en Informatique) est une association dédiée à la recherche et à l'innovation dans le domaine de l'informatique, dont les activités couvrent la conduite de projets de recherche, le développement de solutions technologiques et l'organisation d'activités de formation et de diffusion scientifique.

Dans le cadre de sa mission, l'ARSII développe et anime un réseau de partenaires et d'experts dans les domaines de la R\&I, en lien avec la coopération scientifique entre l'Europe et l'Afrique. La gestion de ce réseau est au cœur de ses activités : identifier des experts, mobiliser des compétences pour des projets et suivre les relations avec les partenaires. Ce besoin de gestion justifie le projet mené durant ce stage : l'association a besoin d'un outil dédié, adapté à la spécificité de son domaine, pour organiser, enrichir et exploiter sa base de contacts.

\section{Étude de l'existant}

Avant la réalisation de la plateforme, la gestion des contacts de l'ARSII reposait soit sur des pratiques essentiellement manuelles, soit sur le recours à des solutions CRM du marché, deux options dont les limites et les coûts étaient devenus structurants.

\subsection{Limites de la gestion manuelle}

La gestion manuelle, majoritairement fondée sur des fichiers tableurs (Excel), présentait plusieurs limites, résumées dans le tableau~\ref{tab:limites-manuelles}.

\begin{table}[H]
\centering
\small
\begin{tabularx}{\textwidth}{|l|X|}
\hline
\textbf{Limite} & \textbf{Conséquence} \\
\hline
Dispersion des données & Contacts répartis dans plusieurs fichiers gérés individuellement, sans référentiel commun : vision globale impossible, doublons, incohérences et données obsolètes \\
\hline
Absence de traçabilité & Aucune trace fiable des ajouts, modifications ou suppressions, ni de l'origine des informations (saisie manuelle, import, numérisation) \\
\hline
Exploitation limitée & Recherche fine et segmentation selon des critères pertinents (pays, affiliation, domaine de compétence) impossibles sans manipulations manuelles fastidieuses \\
\hline
Saisie chronophage & Cartes de visite collectées lors d'événements ressaisies à la main, une opération longue et source d'erreurs \\
\hline
\end{tabularx}
\caption{Limites de la gestion manuelle des contacts}
\label{tab:limites-manuelles}
\end{table}

\subsection{Les solutions CRM du marché}

Pour combler ces manques, l'association pourrait aussi se tourner vers des solutions CRM (Customer Relationship Management) matures, telles que Salesforce, HubSpot ou Zoho CRM. Malgré leurs atouts (gestion centralisée, automatisation, reporting), elles restent inadaptées aux besoins spécifiques de l'ARSII, comme le montre la comparaison du tableau~\ref{tab:crm}.

\begin{table}[H]
\centering
\small
\begin{tabularx}{\textwidth}{|l|X|X|}
\hline
\textbf{Critère} & \textbf{Solutions CRM du marché} & \textbf{Ce qu'une solution dédiée devrait offrir} \\
\hline
Coût & Abonnement souvent payant par utilisateur, budget non négligeable pour une association & Un coût maîtrisé et pérenne, sans abonnement récurrent \\
\hline
Adaptation au domaine & Outils génériques, non pensés pour le contexte R\&I Europe--Afrique & Une adaptation au contexte R\&I Europe--Afrique \\
\hline
Numérisation des cartes & Occasionnelle ou absente, souvent en option & Une numérisation \ac{OCR} intégrée à la saisie \\
\hline
Assistant conversationnel & Souvent payant ou indisponible & Un assistant IA en langage naturel intégré \\
\hline
Contrôle des données & Données hébergées chez un prestataire externe & Des données maîtrisées par l'association \\
\hline
Centralisation et traçabilité & Centralisation possible, mais dépendante de l'abonnement & Une centralisation et une traçabilité complètes \\
\hline
\end{tabularx}
\caption{Comparaison des solutions CRM du marché avec les exigences d'une solution dédiée}
\label{tab:crm}
\end{table}

Au-delà du coût, ces solutions restent insuffisantes sur les points clés du besoin : inadaptées à la spécificité du domaine R\&I Europe--Afrique, elles placent en outre les données du réseau entre les mains d'un prestataire externe.

Il en ressort un besoin non couvert : un outil unique, centralisé et dédié pour structurer, enrichir et exploiter le réseau de contacts R\&I. C'est de ce constat que découle la problématique exposée ci-après.

\section{Problématique et objectifs}

\subsection{Problématique}

Le constat dressé par l'étude de l'existant soulève la question centrale suivante :

\begin{quotation}
\textit{Comment concevoir et développer une plateforme web permettant à l'ARSII de gérer, d'organiser et d'enrichir sa base de contacts R\&I de manière centralisée, fiable et efficace, tout en automatisant au maximum les tâches de saisie et d'exploitation du réseau ?}
\end{quotation}

\subsection{Objectifs}

À partir de cette problématique, quatre objectifs mesurables ont été fixés, détaillés dans le tableau~\ref{tab:objectifs}.

\begin{table}[H]
\centering
\small
\begin{tabularx}{\textwidth}{|l|X|}
\hline
\textbf{Objectif} & \textbf{Description} \\
\hline
Centraliser & Rassembler l'ensemble des contacts dans une base de données unique, structurée et persistante, offrant une vision globale du réseau \\
\hline
Automatiser la saisie & Permettre l'import de fichiers (CSV, Excel) et la numérisation de cartes de visite par \ac{OCR}, afin de réduire drastiquement la ressaisie manuelle \\
\hline
Organiser et exploiter & Offrir une recherche et un filtrage avancés, une segmentation du réseau par ensembles de filtres sauvegardés et une catégorisation par tags, et un tableau de bord interactif de synthèse \\
\hline
Sécuriser & Garantir l'accès par une authentification à deux rôles (Utilisateur et Administrateur), l'administration des comptes étant réservée à l'administrateur \\
\hline
\end{tabularx}
\caption{Objectifs mesurables du stage}
\label{tab:objectifs}
\end{table}

Les quatre objectifs couvrent ainsi l'ensemble du cycle de vie des données, de l'alimentation de la base à son exploitation et à sa sécurisation.

\section{Méthodologie de gestion de projet}

La conduite du projet a reposé sur une méthodologie inspirée de l'\textbf{Agile/Scrum}. Face à un périmètre large, découpé en plusieurs modules, une approche itérative et incrémentale s'est révélée adaptée : elle a permis de livrer des fonctionnalités testables à un rythme régulier et de faire valider progressivement les choix de conception et de réalisation par l'encadrant.

Concrètement, le travail a été organisé en sprints thématiques, chacun dédié à un ensemble cohérent de fonctionnalités :

\begin{itemize}[leftmargin=*,itemsep=3pt]
  \item la préparation de l'environnement et du socle technique.
  \item le module de gestion des contacts et d'import.
  \item le module de recherche et de segmentation.
  \item l'assistant conversationnel IA.
  \item la numérisation OCR.
  \item l'administration et la sécurisation.
\end{itemize}
